A rendszer két fő részre bontható fel: háttérfolyamatok és megjelenítés. Mivel én a megjelenítés résszel foglalkozom, amely magában foglalja a felhasználói weboldalt, az adminisztrátori weboldalt és a mobil applikációt, ezért a következőkben csak ezekre a részekre vonatkozó követelmények specifikációját fogom részletezni. A háttérfolyamatok követelményei Zsigmond Attila kollégám dolgozatában lesznek ismertetve \cite{bus4userver}.

\section{Felhasználói követelmények}
A felhasználói követelmények azokat a funkciókat írják le, amelyre a felhasználó szemszögéből szükség van. A szoftver általános felhasználói követelményei a következők:

\begin{itemize}
  \item A felhasználó gyorsan és egyszerűen kell tudjon navigálni az oldalak között.
  \item Minden felhasználói bevitel esetén a rendszernek biztosítania kell a kényelmes használatot, valamint a megadott adatok helyességéről visszajelzést kell biztosítania beszédes hibaüzenetek megjelenítése által.
  \item A felhasználói felületnek reszponzívnak kell lennie, azaz minden eszközön jól használhatónak kell lennie, legyen az mobiltelefon, tablet vagy számítógép.
  \item A felhasználói felületnek egységesnek kell lennie, azaz minden oldalon ugyanazok a színek és felületi elemek jelennek meg, hogy a felhasználók ne érezzék zavarónak az oldalak közötti váltást.
\end{itemize}

A frontend három különböző felületet tartalmaz: mobilalkalmazás, felhasználói- és adminisztrátori weboldal, ezek pedig különböző felhasználói követelményekkel rendelkezhetnek. A mobilapplikáció ugyanúgy működik, mint a felhasználóknak szánt weboldal, így ezeknek a követelményeit együtt fogom tárgyalni felhasználói felület néven. Ezzel szemben az adminisztrátori weboldal különálló egység, amelynek saját követelményei vannak.

\subsection{Felhasználói felület}
A felhasználói felület funkcióinak nagy része bejelentkezés nélkül használható. A fiókot igénylő funkciók a jegyvásárlás és a jegyek listázása. Az ezek eléréséhez szükséges regisztráció és bejelentkezés bármely oldalról elérhető. A felhasználói felület funkciói oldalanként csoportosítva a következők:

\begin{itemize}
  \item Főoldal: járatokra való rákeresés dátum, idő, kezdő- és célállomás alapján, találatok listázása indulási idő szerinti növekvő sorrendben, jegyvásárlás (bejelentkezés szükséges)
  \item Menetrendek: buszok és megállók órarendjeinek megtekintése, járatok útvonalának megjelenítése térképen
  \item Megállók: megállók listázása térképen, szűrés település vagy buszjárat alapján
  \item Élő térkép: buszok valós idejű pozícióinak megtekintése térképen
  \item Jegyek (bejelentkezés szükséges): megvásárolt jegyek listázása QR-kódokkal
\end{itemize}

\subsection{Adminisztrátori felület}
Az adminisztrátori felület kizárólag bejelentkezéssel elérhető, új fiókokat pedig csak a
cégek főnökei tudnak létrehozni, mindenik a saját busztársaságához. A társult vállalatok főnökei a partnerség kezdetekor megkapják az admin jogú fiókjaikat. Az adminfelület funkciói oldalanként csoportosítva:

\begin{itemize}
  \item Főoldal: statisztikák az eladott jegyek számáról, a megtett kilométerről és a bevételről, időintervalumok szerint csoportosítva, értesítések megjelenítése
  \item Megállók kezelése: megállók listázása, új megálló hozzáadása térkép segítségével, meglévő megálló szerkesztése, megállók törlése
  \item Járatok kezelése: buszjáratok listázása, új járat hozzáadása, meglévő járat szerkesztése: megállók hozzáadása és eltávolítása, járatok törlése
  \item Menetrendek kezelése: járatok indulási időpontjainak és jegyárainak megtekintése, szerkesztése és törlése
  \item Buszok kezelése: buszok listázása, új busz hozzáadása, meglévő busz szerkesztése: útadó, biztosítás és műszaki vizsga időpontok megadása, buszok törlése
  \item Alkalmazottak kezelése: alkalmazotti fiókok listázása, új alkalmazott hozzáadása, meglévő alkalmazott szerkesztése, alkalmazottak törlése
\end{itemize}

Belépés előtt bármely oldalra rákeresve, a rendszer átirányítja a felhasználót a bejelentkező felületre. Regisztrációs lehetőség nincs, így meg lehet akadályozni az adminfelülethez való illetéktelen hozzáférést.

A teljes rendszer funkcionalitását leíró használati eset diagram a \ref{fig:use-case}-es ábrán látható.

\begin{figure}[H]
  \includegraphics[height=\textheight]{figures/use-case.png}
  \caption{Use-case diagram a felhasználók és az adminisztrátorok lehetőségeiről}
  \label{fig:use-case}
\end{figure}

\section{Rendszerkövetelmények}
\subsection{Funkcionális követelmények}

A funkcionális követelmények a felületen megjelenő funkciók működését írják le. Ezek foglalják össze, hogy adott bemenetre milyen kimenetet kell szolgáltatnia a rendszernek. A Bus4U esetében ezeket is két csoportba sorolhatjuk: felhasználói és adminisztrátori követelmények.

A felhasználói weboldalon és a mobilalkalmazásban bejelentkezés nélkül hozzá lehet férni a főoldalhoz, a menetrendekhez, a megállókhoz és az élő térképhez. Az egyetlen oldal, ami nem jeleníthető meg azonosítás nélkül az a jegyek oldal, innen a rendszernek át kell irányítania a bejelentkezési felületre. Itt biztosítani kell a regisztráció lehetőségét is, amennyiben a felhasználó még nem rendelkezik fiókkal. A regisztráció során a rendszer le kell ellenőrizzen minden megadott adatot, még az email tulajdonjogát is, azáltal hogy egy ellenőrző kódot küld a megadott email címre, amelyet be kell írni a felületen megjelenő szövegmezőbe. Sikeres regisztráció vagy bejelentkezés után a rendszernek minden adatot biztonságosan el kell mentenie, a jelszavat hash-elt formában és át kell irányítania a felhasználót a főoldalra.

A főoldalon lévő szövegmezőkben megadott adatok alapján a rendszernek le kell kérnie a szerverről néhány lehetséges járatot, amelyek indulnak az adott településről, adott időben, a megadott célállomásra. A megjelenített találatokhoz tartoznia kell egy jegyvásárlási lehetőségnek is, amelyben meg kell adni a jegyek számát és véglegesíteni kell a vásárlást. A jegyvásárlás is bejelentkezést igényel, így az azonosítatlan felhasználót a rendszer vásárlás előtt átirányítja a bejelentkezéshez. A vásárlás csak a Stripe oldalon történő kifizetés után véglegesítődhet, ezt követi a jegy kódjának kigenerálása és elmentése az adatbázisban. Ezek után, a Jegyek oldalon, a felhasználóknak lehetőségük van a jegyek megtekintésére is, amelyeket korábban vásároltak, és amelyek QR kódos formában jelennek meg a buszon való könnyebb beolvasás végett. A menetrendek oldalon a felhasználóknak lehetőségük van a buszok és megállók közötti választásra lenyíló listákból, amelyet követően a rendszer kilistázza az adott járat megállóhoz kötött időpontjait egy táblázatban. Ugyanitt a rendszernek lehetőséget kell biztosítania a buszok útvonalának megjelenítésére térképen. A megállók oldalon a felhasználóknak lehetőségük van az összes megálló megtekintésére térképen, valamint lenyíló listákat használva, a megállók szűrésére település vagy buszjárat alapján is. A térképnek interaktívnak kell lennie, így a megjelölt megállókra való kattintás után meg kell jelennie az adott megálló adatainak egy kis előugró ablakban. Az élő térkép oldalon a felhasználóknak lehetőségük van a buszok valós idejű pozícióinak követésére térképen. A megjelenített buszokat a rendszernek valós időben kell frissítenie, legalább félpercenként, hogy a felhasználók pontos információkat kapjanak a buszok helyzetéről. Az összes térképes oldalon a felhasználóknak lehetőségük van a térkép nagyítására és kicsinyítésére, valamint a térkép mozgatására is, egérrel vagy érintéses gesztusok segítségével. A felhasználónak meg kell tudnia tekinteni a saját pozícióját is a térképeken, amennyiben engedélyezte a GPS-hez való hozzáférést.

Az adminisztrátori weboldalon szükség van azonosításra, anélkül a rendszer minden oldalról átirányít a bejelentkezéshez. Itt a regisztráció nem lehetséges, így külső, admin szerepkörrel nem rendelkező felhasználók nem tudnak hozzáférni az adminisztrátori felülethez. A főoldalon található statisztikák között szerepel az eladott jegyek száma, a megtett kilométerek száma és a bevétel összege, időintervallumok szerint csoportosítva, táblázatos formában. Mivel minden busztársaságnak el kell legyenek különítve az adatai, ezért a szerver mindig a bejelentkezett adminhoz tartozó cég adataiból számolja ki a statisztikákat. A főoldalon a rendszernek meg kell jelenítenie az esetleges értesítéseket is, előugró "Toast" üzenetek segítségével. A megállók oldalon az adminisztrátoroknak lehetőségük van a megállók listázására, az ezekre való egyenkénti kattintás után pedig a szerkesztésükre is. Emellett megállók törlésére és új megállók hozzáadására is lehetőség van egy interaktív térkép segítségével. A térképre kattintva a hely koordinátáit a rendszer beilleszti az űrlapba, ahol meg kell adni még a megálló nevét és a címet is. Utóbbit, amennyiben elérhető, le kell kérnie a rendszernek a térképszolgáltatótól, hogy ezzel is gyorsítsa a létrehozást. A járatok oldalon az adminisztrátoroknak lehetőségük van a buszjáratok kiválasztására egy lenyíló menüből, ezt követően pedig meg kell jelennie a járatszerkesző űrlapnak. Ugyanitt lehetőség lesz járatok hozzáadására és törlésére is. A menetrendek oldalon az adminisztrátoroknak lehetőségük van a járatok indulási időpontjainak és jegyárainak megtekintésére, szerkesztésére és törlésére. Az órarendek bevitelét le kell lehessen bontani napokra, ezeket pedig dinamikusan megjeleníteni, hogy minél kevesebb helyet foglaljon az oldalon. A buszok és alkalmazottak oldalakon a rendszernek biztosítania kell egy egységes felületet, amelyben lehetőség nyílik az elemek listázására, új elemek hozzáadására, meglévő elemek szerkesztésére és ezek törlésére is. A buszoknál fontos megadni az útadó, biztosítás és műszaki vizsga időpontjait, mert a rendszer ezek alapján fog emlékeztetőket megjeleníteni a felületre való belépések alkalmával. Az alkalmazottaknál meg kell adni a nevet, a beosztást és a telefonszámot, valamint az emailt és a jelszót is, amelyet a rendszer hash-elt formában tárol. Utóbbiak az alkalmazott bejelentkezéséhez lesznek szükségesek, így a rendszernek ellenőriznie kell őket.

\subsection{Nem funkcionális követelmények}
\subsubsection{Szervezéssel kapcsolatos követelmények}

\begin{itemize}
  \item Programozási nyelvek: Dart, HTML, CSS, JavaScript
  \item Keretrendszerek: Flutter, Vue.js
  \item IDE: Visual Studio Code
  \item Verziókövetés: Git és GitHub
  \item Webhost: Netlify
  \item Csoportos kommunikáció: Slack
\end{itemize}

\subsubsection{Termékkel kapcsolatos követelmények}

A felhasználói weboldal és az admin felület használatához egy modern, internetezésre alkalmas eszközre van szükség, a térképfunkciókhoz teljes kihasználásához engedélyezni kell a helymeghatározást. Ajánlott a 2018 után kiadott böngészőverziók használata:
\begin{itemize}
  \item Chrome >= 87
  \item Firefox >= 78
  \item Safari >= 14
  \item Edge >= 88
\end{itemize}

Az mobilapplikáció használatához szükség van egy okostelefonra vagy tabletre,
amelyen minimum Android 5.0 vagy iOS 8.0 operációs rendszer fut. Fontos hogy legyen
rajta internet elérés (3G,4G,5G vagy Wi-Fi), legalább 120MB szabad tárhely és 516-1024 MB szabad memória. Itt is ajánlott a helymeghatározás (GPS) engedélyezése, hogy a térképfunkciók helyesen működjenek.
